\begin{ouverture}[geometrie]{13}{Thalès, réciproque et contraposée}{Partie 4 · Espace et géométrie}
{Alignements · Parallèles · Rapports égaux · Réciproque · Contraposée}
\objectif{Écrire les rapports de longueurs correspondantes.}
\objectif{Calculer une longueur dans les deux configurations.}
\objectif{Démontrer que des droites sont, ou ne sont pas, parallèles.}
\prerequis{Proportionnalité, fractions, produits en croix et droite des milieux.}
\end{ouverture}
\pagestyle{otMain}\essentiel\fontsize{11}{15}\selectfont
\essSection{1}{Le théorème de Thalès}
\begin{propriete}[Théorème de Thalès]Si A, M, B sont alignés, A, N, C sont alignés et $(MN)\parallel(BC)$, dans la configuration des triangles emboîtés ou dans celle dite du papillon, alors
\[\frac{AM}{AB}=\frac{AN}{AC}=\frac{MN}{BC}.\]
Les longueurs des côtés correspondants des triangles AMN et ABC sont proportionnelles.\end{propriete}
\begin{savoirfaire}[Associer les côtés]Dans les triangles AMN et ABC, A correspond à A, M à B et N à C. Tous les numérateurs appartiennent au même triangle. Choisir ensuite l’égalité contenant l’inconnue et deux données.\end{savoirfaire}
\begin{exemple}Si $AM=3$, $AB=5$, $AC=7$ et les hypothèses sont vérifiées, $AN/7=3/5$, donc $AN=4{,}2$ dans la même unité.\end{exemple}
\begin{attention}Si M est entre A et B, $MB=AB-AM$. Dans le papillon, où A est entre M et B, $MB=AM+AB$. Lire l’ordre des points avant de calculer.\end{attention}
\begin{exemple}[Une preuve par les aires]Dans le cas des triangles emboîtés, les triangles MNB et MNC ont la même base [MN] et la même hauteur, car $(BC)\parallel(MN)$. Comparer les aires conduit à l’égalité des rapports : la démonstration complète figure en fin de livre.\end{exemple}
\clearpage
\essSection{2}{Réciproque, contraposée et proportionnalité}
\begin{propriete}[Réciproque]Si A, M, B et A, N, C sont alignés dans le même ordre et si $\dfrac{AM}{AB}=\dfrac{AN}{AC}$, alors $(MN)\parallel(BC)$.\end{propriete}
\begin{propriete}[Contraposée]Si A, M, B et A, N, C sont alignés et si $\dfrac{AM}{AB}\ne\dfrac{AN}{AC}$, alors les droites (MN) et (BC) ne sont pas parallèles : sinon le théorème imposerait l’égalité.\end{propriete}
\begin{savoirfaire}[Comparer deux rapports]Comparer exactement, par exemple avec des produits en croix : $3/8$ et $5/12$ diffèrent, car $3\times12=36$ et $5\times8=40$. Un arrondi ne permet jamais de conclure à l’égalité.\end{savoirfaire}
\begin{propriete}[Lien avec la proportionnalité]Dans un repère, des points alignés avec l’origine ont des coordonnées proportionnelles : Thalès s’applique aux triangles formés avec leurs projetés sur l’axe des abscisses.\end{propriete}
\begin{savoirfaire}[Rédiger]Énoncer les alignements et le parallélisme, nommer le théorème, écrire l’égalité des rapports puis calculer. Dans une réciproque ou une contraposée, conclure en nommant les deux droites.\end{savoirfaire}
\clearpage\rubrique{Découverte}
\begin{blocLivre}{ex-13-1}\begin{exercice}[1]{Écrire les rapports}
Dans la figure, M appartient à [AB], N à [AC] et $(MN)\parallel(BC)$. Écrire les rapports égaux.
\label{fig-ex-13-1}\figureGeo{thales-01}
\end{exercice}\end{blocLivre}
\begin{blocLivre}{ex-13-2}\begin{exercice}[1]{Une longueur réduite}
M est sur [AB], N sur [AC] et $(MN)\parallel(BC)$. AM = 3 cm, AB = 5 cm, AC = 10 cm. Calculer AN.
\label{fig-ex-13-2}\figureGeo{thales-02}
\end{exercice}\end{blocLivre}
\begin{blocLivre}{ex-13-3}\begin{exercice}[1]{La petite base}
Dans la même configuration parallèle, AM = 4 cm, AB = 10 cm et BC = 15 cm. Calculer MN.
\label{fig-ex-13-3}\figureGeo{thales-03}
\end{exercice}\end{blocLivre}
\begin{blocLivre}{ex-13-4}\begin{exercice}[1]{Retrouver le grand côté}
M est sur [AB], N sur [AC], $(MN)\parallel(BC)$. AM = 6 cm, AN = 4 cm et AC = 10 cm. Calculer AB.
\label{fig-ex-13-4}\figureGeo{thales-04}
\end{exercice}\end{blocLivre}
\begin{blocLivre}{ex-13-5}\begin{exercice}[1]{De la partie au tout}
M est sur [AB], N sur [AC], $(MN)\parallel(BC)$. AM = 3 cm, MB = 2 cm et AC = 12 cm. Calculer AN.
\label{fig-ex-13-5}\figureGeo{thales-05}
\end{exercice}\end{blocLivre}

\clearpage\rubrique{Entraînement}
\begin{blocLivre}{ex-13-6}\begin{exercice}[2]{Démontrer un parallélisme}
A, M, B sont alignés dans cet ordre ; A, N, C aussi. AM = 4, AB = 10, AN = 6 et AC = 15, en cm. Montrer que (MN) et (BC) sont parallèles.
\label{fig-ex-13-6}\figureGeo{thales-06}
\end{exercice}\end{blocLivre}
\begin{blocLivre}{ex-13-7}\begin{exercice}[2]{Des droites non parallèles}
A, M, B et A, N, C sont alignés dans cet ordre. AM = 3, AB = 8, AN = 5, AC = 12, en cm. Les droites (MN) et (BC) sont-elles parallèles ?
\label{fig-ex-13-7}\figureGeo{thales-07}
\end{exercice}\end{blocLivre}
\begin{blocLivre}{ex-13-8}\begin{exercice}[2]{Le papillon}
M, A, B sont alignés dans cet ordre ; N, A, C aussi. $(MN)\parallel(BC)$, AM = 3 cm, AB = 5 cm et AC = 8 cm. Calculer AN.
\label{fig-ex-13-8}\figureGeo{thales-08}
\end{exercice}\end{blocLivre}
\begin{blocLivre}{ex-13-9}\begin{exercice}[2]{La distance entre deux points}
Dans le papillon précédent, AM = 3 cm et AB = 5 cm. Calculer MB. Si AN = 4,8 cm et AC = 8 cm, calculer NC.
\label{fig-ex-13-9}\figureGeo{thales-09}
\end{exercice}\end{blocLivre}
\begin{blocLivre}{ex-13-10}\begin{exercice}[2]{Une réciproque en papillon}
M, A, B sont alignés dans cet ordre, ainsi que N, A, C. AM = 2,4 cm, AB = 4 cm, AN = 3 cm et AC = 5 cm. Les droites (MN) et (BC) sont-elles parallèles ?
\label{fig-ex-13-10}\figureGeo{thales-10}
\end{exercice}\end{blocLivre}

\clearpage\rubrique{Problèmes}
\begin{blocLivre}{ex-13-11}\begin{exercice}[3]{Deux triangles de même aire}
Dans le triangle ABC, M appartient à [AB], N à [AC] et $(MN)\parallel(BC)$. Expliquer pourquoi les triangles MNB et MNC ont la même aire. En déduire que les triangles AMC et ANB ont la même aire.
\label{fig-ex-13-11}\figureGeo{thales-aires}
\end{exercice}\end{blocLivre}
\begin{blocLivre}{ex-13-12}\begin{exercice}[3]{L’ombre d’un arbre}
Un bâton vertical de 1,5 m est planté de sorte que l’extrémité S de son ombre coïncide avec celle de l’ombre d’un arbre vertical. L’ombre du bâton mesure 2 m, celle de l’arbre 12 m, sur un sol horizontal. Les rayons du soleil sont rectilignes. Calculer la hauteur de l’arbre.
\label{fig-ex-13-12}\figureGeo{thales-ombre}
\end{exercice}\end{blocLivre}
\begin{blocLivre}{ex-13-13}\begin{exercice}[3]{Une équation dans Thalès}
M est sur [AB], N sur [AC], (MN) est parallèle à (BC). AM = 4 cm, MB = $x$ cm, AN = 6 cm et NC = 3 cm. Trouver $x$.
\label{fig-ex-13-13}\figureGeo{thales-13}
\end{exercice}\end{blocLivre}
\begin{blocLivre}{ex-13-14}\begin{exercice}[3]{Calculer une largeur inaccessible}
A, B, C forment un triangle. M est le milieu de [AB] et N le milieu de [AC] ; [BC] traverse une zone inaccessible. On mesure MN = 7,4 m. Justifier le parallélisme de (MN) et (BC), puis calculer BC.
\label{fig-ex-13-14}\figureGeo{thales-milieux}
\end{exercice}\end{blocLivre}
\begin{blocLivre}{ex-13-15}\begin{exercice}[3]{Un contrôle critique}
Dans une configuration emboîtée avec (MN) parallèle à (BC), un élève écrit $AM/AB=NC/AC$. AM = 3 cm, AB = 8 cm, AC = 16 cm. Corriger et calculer NC.
\label{fig-ex-13-15}\figureGeo{thales-15}
\end{exercice}\end{blocLivre}

\clearpage
